% !TeX root = ./T2_protokoll.tex

\section{Einführung}

Das Ziel dieses Versuches ist, mithilfe von Gammastrahlern die Winkelabhängigkeit des Differentiellen Wirkungsquerschnitts bei Compton Streuung zu untersuchen und mit der Theorie zu vergleichen.\\
Dafür muss erst die Kalibrierungkurve und Effizienz unsere Messgeräts mithilfe von Gammaspektroskopie bestimmt werden. Danach werden die Copmtonpeaks berücksichtigt und mithilfe der Geometrie und Eigeschafften des Strahlers schließlich der Wirkungsquerschnitt über die Verküpfung
\begin{equation}
    m = \frac{AI_\gamma}{4\pi r_0^2}\eta \epsilon N_e \frac{\text{d}\sigma}{\text{d}\Omega}\frac{F_D}{r^2}
\end{equation}
bestimmt bei unterschiedlichen Winkeln $\theta$. 
\\
Zusätzlich wird mithilfe vom Escapepeak bei der Gammaspektroskopie und mithilfe einer Auftragung von $1/E_\gamma'-1/E_\gamma$ die Elektronenmasse bestimmt.

